% GENERATED FILE. Edit canonical lessons, then run npm run generate:books.
% canonical-source-hash: fnv1a64:3e186619261f2548
% canonical-lessons: ES-C05-adios, ES-C05-hasta, ES-C05-hasta-limits, ES-C05-repaso-adios

\chapter{Goodbye}
\label{ch:goodbye}
\hlchaptermodality{11}

\begin{chapteropening}
\textbf{By the end of this chapter:} \emph{I can end a conversation, and I know the word every softer goodbye is built from.}

Goodbye — the chapter run end to end.
\end{chapteropening}

\section[adiós]{adiós — \textquotedblleft{}goodbye,\textquotedblright{} or \textquotedblleft{}to God\textquotedblright{}}

\label{lesson:ES-C05-adios}



\begin{quote}
You can greet, introduce yourself, and ask how someone is. Time to leave gracefully. The final, plainest parting is \textbf{adiós} — and, like the English \emph{goodbye}, it's a tiny prayer worn smooth.
\end{quote}

\begin{sounds}
\begin{itemize}
\raggedright
  \item \texttt{diphthong-io} — the \emph{-iós} is one glide, \emph{ee-OHS}: \emph{ah-DYOHS}, stress on the end (that's what the accent on \emph{ó} marks — remember the acute-accent lesson).
\end{itemize}
\end{sounds}

\begin{cousinweb}
\textbf{adiós} is \textbf{a + Dios} — \textquotedblleft{}\textbf{to God}.\textquotedblright{} It's the front half of an old blessing, \emph{a Dios te encomiendo}, \textquotedblleft{}I commend you \textbf{to God}\textquotedblright{} — said to someone setting off, when a journey was dangerous and you left them in God's keeping.

Every Romance sister kept the same farewell-prayer:

\noindent
\begin{tabularx}{\linewidth}{@{}>{\raggedright\arraybackslash}X>{\raggedright\arraybackslash}X>{\raggedright\arraybackslash}X@{}}
\toprule
\textbf{language} & \textbf{goodbye} & \textbf{literally} \\
\midrule
\textbf{Spanish} & \emph{adiós} & \textquotedblleft{}to God\textquotedblright{} \\
French & \emph{adieu} & \textquotedblleft{}to God\textquotedblright{} (\emph{à Dieu}) \\
Italian & \emph{addio} & \textquotedblleft{}to God\textquotedblright{} \\
Portuguese & \emph{adeus} & \textquotedblleft{}to God\textquotedblright{} \\
\bottomrule
\end{tabularx}

And English is doing the \textbf{exact same thing} without telling you: \textbf{goodbye} is a crushed-down \emph{\textquotedblleft{}God be with ye.\textquotedblright{}} So \emph{adiós} and \emph{goodbye} are, at heart, the same sentence — a blessing for the road.

\begin{itemize}
\raggedright
  \item \textbf{Dios} (\textquotedblleft{}God\textquotedblright{}) $\leftarrow$ Latin \textbf{Deus} — the \emph{Deus} of English \emph{deity}, \emph{divine}, \emph{adieu}, and (via Greek \emph{Zeus}, same root) the whole sky-god family.
\end{itemize}
\end{cousinweb}

\begin{grammarlens}[title={Grammar Lens: adiós is a \emph{final} goodbye}]
\emph{adiós} leans toward a \textbf{longer} parting — you won't see them soon, or it's the end of the day. For \textquotedblleft{}see you later / tomorrow,\textquotedblright{} Spanish reaches for \textbf{hasta …}, which the rest of this chapter is about — and \emph{hasta} has a surprising origin.
\end{grammarlens}

\subsection*{Your turn}
Say these aloud:

\begin{itemize}
\raggedright
  \item \textquotedblleft{}adiós\textquotedblright{} — \emph{ah-DYOHS}, stress the end
  \item adiós / adieu / addio / adeus — four \textquotedblleft{}to God\textquotedblright{} farewells
  \item English \textquotedblleft{}goodbye\textquotedblright{} = \textquotedblleft{}God be with ye\textquotedblright{} — the same prayer
\end{itemize}

\begin{tcolorbox}[breakable,colback=teal!4,colframe=teal!35!black,title={Before you move on}]
What two words hide inside \emph{adiós}? (\emph{a} + \emph{Dios} — \textquotedblleft{}to God.\textquotedblright{}) What English goodbye is built the same way? (\emph{Goodbye} $\leftarrow$ \textquotedblleft{}God be with ye.\textquotedblright{}) Is \emph{adiós} for a quick \textquotedblleft{}see you soon\textquotedblright{}? (No — it's the longer/final parting; \emph{hasta …} comes next.) Next: \textbf{hasta}, and a word Spain borrowed from Arabic.
\end{tcolorbox}
\section[hasta]{hasta — \textquotedblleft{}until,\textquotedblright{} and Spain's Arabic centuries}

\label{lesson:ES-C05-hasta}



\begin{quote}
Spanish says \textquotedblleft{}see you later\textquotedblright{} as \textbf{hasta luego} — literally \textquotedblleft{}\textbf{until} later.\textquotedblright{} So the key word is \textbf{hasta}, \textquotedblleft{}until.\textquotedblright{} And \emph{hasta} carries the single most important fact about Spanish's history: \textbf{for nearly 800 years, much of Spain spoke Arabic.}
\end{quote}

\begin{sounds}
\begin{itemize}
\raggedright
  \item \texttt{silent-h} — Spanish \textbf{h is always silent}: \textbf{hasta} = \emph{AHS-tah} (no breath on the \emph{h} at all). Stress the front.
\end{itemize}
\end{sounds}

\begin{cousinweb}
\textbf{hasta} comes from Arabic \textbf{ḥattā}, “until, up to.” From \textbf{711 CE}, Muslim armies ruled much of the Iberian Peninsula — \emph{al-Andalus} — and Arabic was the language of government, science, and poetry there for \textbf{centuries}, until 1492. Spanish absorbed roughly \textbf{four thousand} Arabic words in that time — more of its vocabulary than from any source except Latin itself.

Most of those loans are \textbf{nouns}, and you can spot a huge number by the frozen Arabic article \textbf{al-} (\textquotedblleft{}the\textquotedblright{}) fused onto the front:

\begin{itemize}
\raggedright
  \item \textbf{almohada} (pillow), \textbf{alcázar} (fortress), \textbf{alcalde} (mayor), \textbf{alfombra} (rug), \textbf{aduana} (customs), \textbf{álgebra}, \textbf{alcohol}.
  \item \textbf{azúcar} (sugar), \textbf{aceite} (oil), \textbf{naranja} (orange).
\end{itemize}

But \textbf{hasta} is rare and special: it is a \textbf{function word} — a preposition, part of the grammatical skeleton. Languages borrow nouns easily and grammar words almost never — so an Arabic \emph{preposition} surviving in everyday Spanish shows how deep, how \emph{intimate}, those centuries of contact were. (Its Latin rival \emph{ad} \textquotedblleft{}to\textquotedblright{}/\emph{usque} \textquotedblleft{}up to\textquotedblright{} simply lost.) You say a piece of al-Andalus every time you say goodbye.
\end{cousinweb}

\subsection*{Your turn}
Say these aloud:

\begin{itemize}
\raggedright
  \item \textquotedblleft{}hasta\textquotedblright{} — \emph{AHS-tah}, silent h
  \item \textquotedblleft{}hasta\textquotedblright{} then Arabic \emph{ḥattā} — the same word, 1,300 years apart
\end{itemize}

\begin{tcolorbox}[breakable,colback=teal!4,colframe=teal!35!black,title={Before you move on}]
What language is \emph{hasta} from, and what does it mean? (Arabic \emph{ḥattā} — \textquotedblleft{}until.\textquotedblright{}) Why is it remarkable that Spanish borrowed it? (It's a \emph{function word} / preposition — languages almost never borrow those; it shows how deep the Arabic centuries went.) How do you spot many Arabic nouns in Spanish? (The frozen \emph{al-} article: \emph{almohada, azúcar, álgebra}.)
\end{tcolorbox}
\section[hasta aquí / hasta la noche]{hasta aquí / hasta la noche — point to the limit}

\label{lesson:ES-C05-hasta-limits}



\begin{quote}
You know \textbf{hasta} means “until” or “up to.” Give that small word one spatial endpoint and one temporal endpoint, using only one new word.
\end{quote}

\begin{grammarlens}[title={Grammar Lens: endpoints in space and time}]
\textbf{Aquí} means \textbf{here}. Put it after \textbf{hasta}:

\begin{quote}
\textbf{hasta aquí} — up to here
\end{quote}

For a time endpoint, reuse \textbf{la noche}, “the night,” which you already have:

\begin{quote}
\textbf{hasta la noche} — until the night / until tonight
\end{quote}

The same preposition points to a place or the end of a stretch of time. The next lessons will place new time words after it one at a time.
\end{grammarlens}

\subsection*{Your turn}
Say these aloud:

\begin{itemize}
\raggedright
  \item “aquí” — here
  \item “hasta aquí” — up to here
  \item “hasta la noche” — until tonight
\end{itemize}

\begin{tcolorbox}[breakable,colback=teal!4,colframe=teal!35!black,title={Before you move on}]
What does \textbf{hasta} point to? (\textbf{An endpoint.}) Give one spatial and one temporal example. (\textbf{Hasta aquí; hasta la noche.}) What is the only new lexical item? (\textbf{Aquí}, “here.”)
\end{tcolorbox}
\section[Practice]{Leaving — putting it together}

\label{lesson:ES-C05-repaso-adios}



\begin{quote}
One farewell that is final, and one word that opens all the others. Nothing new in this lesson — only what you already have.
\end{quote}

\subsection*{The exchange}
\begin{quote}
\textbf{Adiós} — meant, and a little final. \textbf{Hasta} — \textquotedblleft{}until,\textquotedblright{} and it is waiting for what comes next.
\end{quote}

\begin{grammarlens}[title={What you've built}]
You can close a conversation, and you have the word every softer goodbye is built from.
\end{grammarlens}

\subsection*{Your turn}
Say these aloud:

\begin{itemize}
\raggedright
  \item the whole exchange, both parts, without stopping
  \item it again, faster, and let it sound like speech rather than recital
\end{itemize}

\emph{Twice through:}

\begin{tcolorbox}[breakable,colback=teal!4,colframe=teal!35!black,title={Before you move on}]
Run it once more from memory. If a word does not come, go back one lesson rather than pushing on — this chapter is short on purpose.
\end{tcolorbox}
